% =============================================================================
% Capítulo 2 — Referencial Teórico
% =============================================================================
\chapter{Referencial Teórico}
\label{cap:referencial-teorico}

Este capítulo apresenta os fundamentos teóricos que sustentam o presente
trabalho. São abordados os conceitos de Máquinas de Estados Finitos e sua
aplicação como técnica de controle de NPCs em jogos digitais, os princípios do
Aprendizado por Reforço com ênfase no algoritmo \textit{Proximal Policy
Optimization}, a plataforma Unity ML-Agents como ambiente de treinamento, o
conceito de Proveniência de Dados e sua aplicação ao rastreamento causal em
sessões de jogo, e a técnica de \textit{Reward Shaping} como mecanismo de
orientação do aprendizado.

% -----------------------------------------------------------------------------
\section{Máquinas de Estados Finitos em Jogos}
\label{sec:fsm}

Uma Máquina de Estados Finitos (FSM --- \textit{Finite State Machine}) é um
modelo computacional composto por um conjunto finito de estados, um conjunto de
transições entre esses estados e um estado inicial. Em cada instante, o sistema
encontra-se em exatamente um estado, e a ocorrência de um evento ou o
atendimento de uma condição pré-definida dispara uma transição para outro
estado \cite{millington2019ai}.

No contexto de jogos digitais, FSMs constituem uma das técnicas mais
difundidas para o controle de NPCs. Cada estado representa um comportamento
distinto do personagem --- como perseguir, atacar, recuar ou permanecer em
espera --- e as transições são determinadas por condições do ambiente, como a
distância até o jogador, o nível de vida remanescente ou a expiração de um
temporizador de \textit{cooldown}. A \autoref{fig:fsm-exemplo} ilustra, de
forma conceitual, a estrutura de uma FSM típica para um chefe de jogo 2D.

 \begin{figure}[htb]
     \centering
     \includegraphics[width=0.7\textwidth]{imagens/fsm_boss.png}
     \caption{Exemplo conceitual de FSM para um chefe de jogo 2D, com
              estados \textit{Idle}, \textit{Chase}, \textit{Attack} e
              \textit{Retreat}.}
     \label{fig:fsm-exemplo}
     \fonte{Elaborado pelo autor.}
 \end{figure}

As principais vantagens de FSMs residem em sua simplicidade de implementação,
previsibilidade comportamental e facilidade de depuração. Essas qualidades
tornam-nas particularmente adequadas para contextos de produção, nos quais os
projetistas necessitam de controle total sobre o comportamento do NPC e de
garantias quanto à reprodutibilidade de cenários de teste
\cite{millington2019ai}.

Contudo, a previsibilidade que constitui uma vantagem em termos de engenharia
converte-se em limitação sob a perspectiva do jogador: após poucas interações,
os padrões de comportamento tornam-se evidentes, reduzindo o desafio percebido
e comprometendo a rejogabilidade \cite{yannakakis2018artificial}. Além disso,
a complexidade de manutenção cresce de forma combinatória com o número de
estados e transições, tornando difícil a adição de novos comportamentos sem
introduzir efeitos colaterais indesejados.

Essas limitações motivam a investigação de abordagens complementares, capazes
de gerar comportamentos menos previsíveis e mais adaptáveis às ações do
jogador. Neste trabalho, a FSM é empregada como \textit{baseline}
determinístico, fornecendo um referencial estável e interpretável para
comparação com o agente treinado por Aprendizado por Reforço.

% -----------------------------------------------------------------------------
\section{Aprendizado por Reforço}
\label{sec:aprendizado-reforco}

O Aprendizado por Reforço (\textit{Reinforcement Learning} --- RL) é um
paradigma de aprendizado de máquina no qual um agente aprende a tomar decisões
por meio de interações com um ambiente. Diferentemente do aprendizado
supervisionado, no qual o modelo é treinado com pares de entrada e saída
rotulados, no RL o agente recebe sinais de recompensa numérica que indicam a
qualidade relativa de suas ações, sem instrução explícita sobre qual ação
tomar em cada situação \cite{sutton2018reinforcement}.

\subsection{Formulação como Processo de Decisão de Markov}
\label{subsec:mdp}

O arcabouço formal do RL é tipicamente modelado como um Processo de Decisão de
Markov (MDP --- \textit{Markov Decision Process}), definido pela tupla
$\langle S, A, P, R, \gamma \rangle$, em que:

\begin{itemize}
    \item $S$ é o conjunto de estados do ambiente;
    \item $A$ é o conjunto de ações disponíveis ao agente;
    \item $P(s' | s, a)$ é a função de transição, que define a probabilidade
          de o ambiente transitar para o estado $s'$ dado que o agente executou
          a ação $a$ no estado $s$;
    \item $R(s, a, s')$ é a função de recompensa, que atribui um valor
          numérico à transição de $s$ para $s'$ por meio da ação $a$;
    \item $\gamma \in [0, 1]$ é o fator de desconto, que pondera a importância
          relativa de recompensas futuras em relação a recompensas imediatas.
\end{itemize}

O objetivo do agente é aprender uma política $\pi(a|s)$ --- uma função que
mapeia estados a distribuições de probabilidade sobre ações --- que maximize o
retorno esperado acumulado ao longo do tempo. O retorno $G_t$ a partir do
instante $t$ é definido como:

\begin{equation}
    G_t = \sum_{k=0}^{\infty} \gamma^k \, R_{t+k+1}
    \label{eq:retorno}
\end{equation}

\noindent em que $R_{t+k+1}$ denota a recompensa recebida $k+1$ passos após o
instante $t$ \cite{sutton2018reinforcement}.

\subsection{Métodos Policy Gradient}
\label{subsec:policy-gradient}

Uma família importante de algoritmos de RL são os métodos de gradiente de
política (\textit{Policy Gradient}), nos quais a política $\pi_\theta$ é
parametrizada diretamente --- por exemplo, por uma rede neural com parâmetros
$\theta$ --- e otimizada por gradiente ascendente sobre uma função objetivo
que mede o retorno esperado. A vantagem dessa família de métodos é a
capacidade de lidar naturalmente com espaços de ação discretos e contínuos,
além de produzir políticas estocásticas que favorecem a exploração
\cite{sutton2018reinforcement}.

O gradiente da função objetivo, segundo o Teorema do Gradiente de Política, é
dado por:

\begin{equation}
    \nabla_\theta J(\theta)
    = \mathbb{E}_{\pi_\theta}
      \!\left[\,
        \nabla_\theta \log \pi_\theta(a_t | s_t) \; \hat{A}_t
      \,\right]
    \label{eq:policy-gradient}
\end{equation}

\noindent em que $\hat{A}_t$ é uma estimativa da função de vantagem
(\textit{advantage function}), que mede o quanto a ação $a_t$ é melhor ou pior
do que a ação média sob a política corrente.

\subsection{Proximal Policy Optimization (PPO)}
\label{subsec:ppo}

O \textit{Proximal Policy Optimization} (PPO), proposto por
\citeonline{schulman2017ppo}, é um algoritmo de gradiente de política que
endereça um problema recorrente em métodos anteriores: a instabilidade causada
por atualizações excessivamente grandes nos parâmetros da política. Algoritmos
como o \textit{Trust Region Policy Optimization} (TRPO) abordaram essa questão
por meio de restrições de segunda ordem, mas com custo computacional elevado.
O PPO propõe uma solução mais simples, baseada em uma função objetivo
\textit{clipped} que limita a magnitude das atualizações sem exigir cálculos
de segunda ordem.

A função objetivo do PPO é definida como:

\begin{equation}
    L^{CLIP}(\theta)
    = \mathbb{E}_t \!\left[\,
        \min\!\Big(
            r_t(\theta) \, \hat{A}_t, \;
            \text{clip}\big(r_t(\theta),\; 1 - \epsilon,\; 1 + \epsilon\big)
            \, \hat{A}_t
        \Big)
    \,\right]
    \label{eq:ppo-clip}
\end{equation}

\noindent em que $r_t(\theta) = \frac{\pi_\theta(a_t | s_t)}
{\pi_{\theta_{\text{old}}}(a_t | s_t)}$ é a razão entre a probabilidade da
ação sob a política atualizada e sob a política anterior, e
$\epsilon$ é um hiperparâmetro que controla o tamanho máximo permitido
para a atualização. Valores típicos de $\epsilon$ situam-se em torno de
$0{,}2$ \cite{schulman2017ppo}.

O mecanismo de \textit{clipping} opera da seguinte forma: quando a razão
$r_t(\theta)$ ultrapassa o intervalo $[1 - \epsilon, \; 1 + \epsilon]$, o
gradiente é atenuado, impedindo que a política se afaste excessivamente da
versão anterior em uma única iteração. Esse mecanismo garante atualizações
conservadoras, resultando em maior estabilidade durante o treinamento.

Do ponto de vista prático, o PPO apresenta características que o tornam
particularmente adequado para aplicações em jogos: estabilidade de treinamento
com ajuste mínimo de hiperparâmetros, compatibilidade com espaços de ação
discretos --- comuns em jogos 2D --- e eficiência amostral razoável em
ambientes simulados. Essas qualidades justificam sua adoção como algoritmo
padrão no Unity ML-Agents \cite{juliani2020unity}.

% -----------------------------------------------------------------------------
\section{Unity ML-Agents}
\label{sec:ml-agents}

O Unity ML-Agents \textit{Toolkit} é uma plataforma de código aberto
desenvolvida pela Unity Technologies que permite o treinamento de agentes
inteligentes em ambientes construídos com o motor de jogo Unity
\cite{juliani2020unity}. A plataforma integra o ambiente de simulação do Unity
com bibliotecas de aprendizado de máquina escritas em Python, notadamente o
PyTorch, estabelecendo um canal de comunicação bidirecional entre o motor de
jogo e o processo de treinamento.

\subsection{Arquitetura da Plataforma}
\label{subsec:ml-agents-arquitetura}

A arquitetura do ML-Agents organiza-se em três camadas principais:

\begin{enumerate}
    \item \textbf{Ambiente Unity}: contém a cena de simulação, os objetos de
          jogo e os scripts de comportamento dos agentes. Cada agente é
          implementado como um componente que herda da classe \texttt{Agent},
          sobrescrevendo métodos para coletar observações, executar ações e
          receber recompensas.

    \item \textbf{Comunicador}: protocolo de comunicação (\textit{side
          channel}) que transmite observações do Unity para o processo Python e
          retorna decisões de ação ao ambiente.

    \item \textbf{Treinador Python}: processo externo que executa o algoritmo
          de aprendizado (PPO, SAC, entre outros), atualiza os parâmetros da
          rede neural do agente e registra métricas de treinamento via
          TensorBoard.
\end{enumerate}

\subsection{Ciclo de Treinamento}
\label{subsec:ml-agents-ciclo}

O ciclo de treinamento no ML-Agents segue a dinâmica padrão do RL. A cada
passo de decisão (\textit{decision step}), o agente coleta observações do
ambiente --- como posições, velocidades e estados internos --- e as envia ao
treinador Python. O treinador seleciona uma ação de acordo com a política
corrente, e o agente a executa no ambiente Unity. O ambiente retorna uma
recompensa e o novo estado, fechando o ciclo.

Ao término do treinamento, o modelo resultante é exportado no formato ONNX
(\textit{Open Neural Network Exchange}) \cite{onnx2019}, um formato
interoperável que permite a execução de inferência diretamente no Unity, sem
a necessidade de manter o processo Python ativo. Essa característica é
fundamental para a integração do agente treinado em um jogo jogável, pois
permite que o modelo opere em modo de inferência local com desempenho
compatível com tempo real.

\subsection{Modos de Operação}
\label{subsec:ml-agents-modos}

O ML-Agents suporta três modos de operação relevantes para este trabalho:

\begin{itemize}
    \item \textbf{Treinamento}: o agente está conectado ao treinador Python e
          atualiza seus parâmetros a cada episódio;
    \item \textbf{Inferência} (\textit{InferenceOnly}): o agente utiliza um
          modelo ONNX previamente treinado para tomar decisões, sem conexão
          com o treinador;
    \item \textbf{Heurístico}: o agente é controlado por lógica programada
          manualmente, útil para coleta de demonstrações ou testes.
\end{itemize}

No presente trabalho, o treinamento foi realizado com PPO via treinador
Python, e o modelo final foi integrado ao Unity em modo
\textit{InferenceOnly} para a etapa de avaliação quantitativa.

% -----------------------------------------------------------------------------
\section{Proveniência de Dados}
\label{sec:proveniencia}

Proveniência de dados (\textit{Data Provenance}) refere-se ao registro
sistemático da origem, das transformações e do histórico de um dado ou
artefato ao longo de seu ciclo de vida. Em sua acepção mais geral, a
proveniência responde às perguntas ``de onde veio este dado?'' e ``por que ele
possui este valor?'', preservando as relações de causa e efeito entre eventos
e entidades em um fluxo de processamento \cite{provenance_base}.

\subsection{Proveniência versus Logs Convencionais}
\label{subsec:proveniencia-vs-logs}

É importante distinguir proveniência de dados de registros de log
convencionais. Enquanto um log tradicional armazena eventos sequenciais com
marcações temporais --- essencialmente respondendo ``o que aconteceu e
quando'' ---, um sistema de proveniência registra, adicionalmente, as relações
causais entre esses eventos. Cada evento possui referências explícitas ao
evento que o originou, formando uma estrutura de grafo dirigido acíclico (DAG)
que representa a cadeia causal da sessão.

Essa distinção é relevante porque a análise de logs planos permite apenas
correlações temporais (``o evento B ocorreu após o evento A''), ao passo que
a análise de proveniência permite inferências causais (``o evento B ocorreu
\textit{em consequência} do evento A''). Essa semântica adicional é
precisamente o que torna a proveniência útil como fonte de informação para
orientar processos de aprendizado.

\subsection{Proveniência em Jogos Digitais}
\label{subsec:proveniencia-jogos}

A aplicação de proveniência ao domínio de jogos digitais foi investigada por
\citeonline{provenance_base}, que propuseram o uso de grafos de proveniência
para a visualização e análise de sessões de jogo. Nessa abordagem, cada ação
relevante do \textit{gameplay} --- como ataques, movimentos, danos recebidos e
mortes --- é registrada como um nó do grafo, e as relações causais entre ações
e consequências são representadas como arestas direcionadas.

\citeonline{kohwalter2020provenance} ampliaram essa investigação com um
levantamento sobre aplicações de proveniência em jogos, demonstrando que os
grafos resultantes podem ser utilizados para análise \textit{post-mortem} de
sessões, identificação de padrões de \textit{gameplay}, detecção de
estratégias emergentes e depuração de \textit{game design}.

A estrutura típica de um grafo de proveniência em uma sessão de jogo de
combate segue um encadeamento causal como o ilustrado a seguir:

\begin{quote}
    \texttt{SessionStart} $\rightarrow$ \texttt{BossAttack} $\rightarrow$
    \texttt{PlayerDamageTaken} $\rightarrow$ \texttt{PlayerAttack}
    $\rightarrow$ \texttt{BossDamageTaken} $\rightarrow$
    \texttt{BossDeath} $\rightarrow$ \texttt{SessionEnd}
\end{quote}

Nessa cadeia, cada evento possui uma referência ao evento que o precedeu
causalmente. A relação entre \texttt{PlayerAttack} e
\texttt{BossDamageTaken}, por exemplo, não é meramente temporal: ela indica
que o dano no chefe foi \textit{causado} pelo ataque do jogador. Essa
distinção é fundamental para a extração de padrões comportamentais
significativos.

No presente trabalho, a proveniência é empregada em dois momentos distintos:
como dado bruto exportado em formato JSON para análise e documentação das
sessões de jogo; e como fonte indireta de informação para o treinamento do
agente, por meio da extração de sequências de ações associadas a desfechos
vitoriosos.

% -----------------------------------------------------------------------------
\section{Modelagem de Recompensa (\textit{Reward Shaping})}
\label{sec:reward-shaping}

A definição de funções de recompensa é uma das etapas mais delicadas do
treinamento por Aprendizado por Reforço. Uma função de recompensa mal
projetada pode levar a comportamentos degenerados: o agente pode encontrar
formas de maximizar a recompensa que não correspondem ao objetivo desejado
pelo projetista --- fenômeno conhecido como \textit{reward hacking}
\cite{sutton2018reinforcement}.

\subsection{Conceito e Fundamentação Teórica}
\label{subsec:reward-shaping-conceito}

\textit{Reward Shaping} é uma técnica que consiste na adição de uma
recompensa auxiliar $F(s, a, s')$ à recompensa original do ambiente $R(s, a,
s')$, com o objetivo de fornecer sinais intermediários que orientem o agente
em direção ao comportamento desejado sem alterar a política ótima do problema
original \cite{ng1999reward}.

\citeonline{ng1999reward} demonstraram formalmente que, quando a recompensa
auxiliar $F$ assume a forma de uma função potencial baseada em estados:

\begin{equation}
    F(s, a, s') = \gamma \, \Phi(s') - \Phi(s)
    \label{eq:pbrs}
\end{equation}

\noindent em que $\Phi: S \rightarrow \mathbb{R}$ é uma função potencial e
$\gamma$ é o fator de desconto, a política ótima do MDP original é preservada.
Essa propriedade, conhecida como invariância de política sob \textit{Potential-
Based Reward Shaping} (PBRS), garante que o \textit{Reward Shaping} não
introduz vícios permanentes no aprendizado: ele apenas acelera a convergência
em direção à política ótima, sem desviá-la.

\subsection{Reward Shaping versus Imitação Direta}
\label{subsec:shaping-vs-imitacao}

É relevante distinguir \textit{Reward Shaping} de técnicas de aprendizagem
por imitação, como GAIL \cite{ho2016gail}. Na imitação direta, o agente
aprende a reproduzir trajetórias demonstradas por um especialista, utilizando
um discriminador adversário que diferencia o comportamento do agente do
comportamento demonstrado. Embora eficaz, essa abordagem vincula o agente ao
comportamento observado, limitando sua capacidade de exploração e dificultando
a emergência de estratégias originais.

O \textit{Reward Shaping}, por outro lado, atua como um sinal orientador que
não substitui a recompensa do ambiente. O agente continua explorando
livremente e otimizando a recompensa total; a componente adicional apenas
indica que certos padrões de comportamento são preferíveis, sem impor uma
trajetória específica. Essa distinção é particularmente relevante no contexto
do presente trabalho, no qual a intenção não é que o chefe \textit{copie} o
comportamento de um jogador humano, mas sim que receba orientação de padrões
causais extraídos de sessões bem-sucedidas.

\subsection{Aplicação ao Contexto do Trabalho}
\label{subsec:shaping-aplicacao}

No pipeline proposto neste trabalho, o \textit{Reward Shaping} é alimentado
por sequências de ações extraídas de grafos de proveniência. Essas sequências
representam padrões de comportamento que precederam eventos favoráveis ---
especificamente, dano causado no chefe em sessões vitoriosas. Quando o agente
PPO, durante o treinamento, executa uma sequência de ações que coincide com
uma das sequências extraídas, ele recebe uma recompensa adicional.

Esse mecanismo combina as vantagens das duas áreas: a proveniência fornece
informação causal interpretável sobre o que constituiu comportamento bem-
sucedido, e o \textit{Reward Shaping} converte essa informação em um sinal
de treinamento compatível com o framework de RL, sem sacrificar a autonomia
exploratória do agente.

% -----------------------------------------------------------------------------
\section{Considerações sobre o Referencial}
\label{sec:consideracoes-referencial}

Os conceitos apresentados neste capítulo fornecem a base teórica para a
compreensão do pipeline proposto. A FSM oferece o \textit{baseline}
determinístico necessário para uma comparação objetiva. O Aprendizado por
Reforço, por meio do PPO, fornece o algoritmo de treinamento do agente. O
Unity ML-Agents viabiliza a execução desse treinamento diretamente no motor
de jogo. A Proveniência de Dados possibilita o registro e a análise causal
de sessões de jogo. E o \textit{Reward Shaping} fornece o mecanismo teórico
para converter padrões causais em sinais de recompensa auxiliares, preservando
a política ótima.

A integração dessas cinco áreas constitui o diferencial metodológico do
presente trabalho, cujos detalhes de implementação e aplicação são
apresentados nos capítulos subsequentes.
